% GENERATED FILE. Edit canonical lessons, then run npm run generate:books.
% canonical-source-hash: fnv1a64:7ccf676dad6a8edc
% canonical-lessons: IT-C285-dichiarazione, IT-C285-portavoce, IT-C285-conferenza-stampa, IT-C285-comunicato, IT-C285-smentita

\chapter{Le parole ufficiali}
\label{ch:it-le-parole-ufficiali}
\hlchaptermodality{285}

\begin{chapteropening}
\textbf{By the end of this chapter:} \emph{I can follow official words: a statement, a spokesperson, a press conference, a press release and a denial.}

Complete the last lesson of chapter 285: I can follow official words: a statement, a spokesperson, a press conference, a press release and a denial.
\end{chapteropening}

\section[la dichiarazione]{la dichiarazione — the statement, the declaration}

\label{lesson:IT-C285-dichiarazione}



\begin{quote}
Before the new one: say the Italian for to testify, then the Italian for to report.
\end{quote}

\subsection*{You'll want to know: la dichiarazione}
\textbf{la dichiarazione} — \textquotedblleft{}the statement, the declaration\textquotedblright{}.

\begin{quote}
Il sindaco ha fatto una breve dichiarazione.
\end{quote}

\begin{quote}
\emph{The mayor made a short statement.}
\end{quote}

\emph{Fare una dichiarazione} — to make a statement.

\begin{cousinweb}
\emph{dichiarazione} is built on \emph{dichiarare} \textquotedblleft{}to declare\textquotedblright{}, inherited from Latin \emph{dēclārāre} \textquotedblleft{}to make clear\textquotedblright{} (EtymDB), from \emph{clārus} \textquotedblleft{}clear\textquotedblright{}.
\end{cousinweb}

\begin{grammarlens}[title={What you've built}]
One more word for telling, explaining and describing at B1.
\end{grammarlens}

\subsection*{Your turn}
\begin{itemize}
\raggedright
  \item \emph{Say it:} \emph{la dichiarazione}
  \item \emph{Say it:} \emph{la dichiarazione}, once more
  \item \emph{Say it:} a sentence of your own with \emph{la dichiarazione}
  \item \emph{Recall:} say the Italian for to testify, then the Italian for to report, then say \emph{la dichiarazione} again
\end{itemize}

\begin{tcolorbox}[breakable,colback=teal!4,colframe=teal!35!black,title={Before you move on}]
What does la dichiarazione mean? (\textquotedblleft{}The statement, the declaration\textquotedblright{}.) Say it once more.
\end{tcolorbox}
\section[il portavoce]{il portavoce — the spokesperson}

\label{lesson:IT-C285-portavoce}



\begin{quote}
Before the new one: say the Italian for to report, then the Italian for the statement.
\end{quote}

\subsection*{You'll want to know: il portavoce}
\textbf{il portavoce} — \textquotedblleft{}the spokesperson\textquotedblright{}.

\begin{quote}
Il portavoce non ha voluto rispondere.
\end{quote}

\begin{quote}
\emph{The spokesperson would not answer.}
\end{quote}

One form for a man or a woman; plural unchanged: \emph{i portavoce}.

\begin{cousinweb}
\emph{portavoce} joins \emph{porta}, from \emph{portare} \textquotedblleft{}to carry\textquotedblright{}, and \emph{voce} \textquotedblleft{}voice\textquotedblright{}, from Latin \emph{vōx} (Lewis \& Short): the one who carries the voice.
\end{cousinweb}

\begin{grammarlens}[title={What you've built}]
One more word for telling, explaining and describing at B1.
\end{grammarlens}

\subsection*{Your turn}
\begin{itemize}
\raggedright
  \item \emph{Say it:} \emph{il portavoce}
  \item \emph{Say it:} \emph{il portavoce}, once more
  \item \emph{Say it:} a sentence of your own with \emph{il portavoce}
  \item \emph{Recall:} say the Italian for to report, then the Italian for the statement, then say \emph{il portavoce} again
\end{itemize}

\begin{tcolorbox}[breakable,colback=teal!4,colframe=teal!35!black,title={Before you move on}]
What does il portavoce mean? (\textquotedblleft{}The spokesperson\textquotedblright{}.) Say it once more.
\end{tcolorbox}
\section[la conferenza stampa]{la conferenza stampa — the press conference}

\label{lesson:IT-C285-conferenza-stampa}



\begin{quote}
Before the new one: say the Italian for the statement, then the Italian for the spokesperson.
\end{quote}

\subsection*{You'll want to know: la conferenza stampa}
\textbf{la conferenza stampa} — \textquotedblleft{}the press conference\textquotedblright{}.

\begin{quote}
Alle undici c'è una conferenza stampa.
\end{quote}

\begin{quote}
\emph{There is a press conference at eleven.}
\end{quote}

\emph{La stampa} is the press; \emph{una conferenza} on its own is a lecture or a talk.

\begin{cousinweb}
\emph{conferenza} is a learned word usually traced to Medieval Latin \emph{cōnferentia}, built on \emph{cōnferre} \textquotedblleft{}to bring together\textquotedblright{} (Lewis \& Short); \emph{stampa} \textquotedblleft{}the press\textquotedblright{} is built on \emph{stampare} \textquotedblleft{}to print\textquotedblright{}, which EtymDB traces to a Frankish verb.
\end{cousinweb}

\begin{grammarlens}[title={What you've built}]
One more word for telling, explaining and describing at B1.
\end{grammarlens}

\subsection*{Your turn}
\begin{itemize}
\raggedright
  \item \emph{Say it:} \emph{la conferenza stampa}
  \item \emph{Say it:} \emph{la conferenza stampa}, once more
  \item \emph{Say it:} a sentence of your own with \emph{la conferenza stampa}
  \item \emph{Recall:} say the Italian for the statement, then the Italian for the spokesperson, then say \emph{la conferenza stampa} again
\end{itemize}

\begin{tcolorbox}[breakable,colback=teal!4,colframe=teal!35!black,title={Before you move on}]
What does la conferenza stampa mean? (\textquotedblleft{}The press conference\textquotedblright{}.) Say it once more.
\end{tcolorbox}
\section[il comunicato]{il comunicato — the press release, the official bulletin}

\label{lesson:IT-C285-comunicato}



\begin{quote}
Before the new one: say the Italian for the spokesperson, then the Italian for the press conference.
\end{quote}

\subsection*{You'll want to know: il comunicato}
\textbf{il comunicato} — \textquotedblleft{}the press release, the official bulletin\textquotedblright{}.

\begin{quote}
L'azienda ha mandato un comunicato ai giornali.
\end{quote}

\begin{quote}
\emph{The company sent a press release to the papers.}
\end{quote}

\emph{Un comunicato stampa} — a press release.

\begin{cousinweb}
\emph{comunicato} is the past participle of \emph{comunicare} \textquotedblleft{}to make known\textquotedblright{}, a learned borrowing from Latin \emph{commūnicāre} \textquotedblleft{}to share\textquotedblright{}, built on \emph{commūnis} \textquotedblleft{}common\textquotedblright{} (Lewis \& Short).
\end{cousinweb}

\begin{grammarlens}[title={What you've built}]
One more word for telling, explaining and describing at B1.
\end{grammarlens}

\subsection*{Your turn}
\begin{itemize}
\raggedright
  \item \emph{Say it:} \emph{il comunicato}
  \item \emph{Say it:} \emph{il comunicato}, once more
  \item \emph{Say it:} a sentence of your own with \emph{il comunicato}
  \item \emph{Recall:} say the Italian for the spokesperson, then the Italian for the press conference, then say \emph{il comunicato} again
\end{itemize}

\begin{tcolorbox}[breakable,colback=teal!4,colframe=teal!35!black,title={Before you move on}]
What does il comunicato mean? (\textquotedblleft{}The press release, the official bulletin\textquotedblright{}.) Say it once more.
\end{tcolorbox}
\section[la smentita]{la smentita — the denial (of a report)}

\label{lesson:IT-C285-smentita}



\begin{quote}
Before the new one: say the Italian for the press conference, then the Italian for the press release.
\end{quote}

\subsection*{You'll want to know: la smentita}
\textbf{la smentita} — \textquotedblleft{}the denial (of a report)\textquotedblright{}.

\begin{quote}
Il giorno dopo è arrivata la smentita.
\end{quote}

\begin{quote}
\emph{The next day the denial came.}
\end{quote}

Official and journalistic: a statement that a story is not true.

\begin{cousinweb}
\emph{smentita} is built on \emph{smentire} \textquotedblleft{}to give the lie to\textquotedblright{}, from \emph{s-} and \emph{mentire} \textquotedblleft{}to lie\textquotedblright{}, which is inherited from Latin \emph{mentīrī} (Lewis \& Short).
\end{cousinweb}

\begin{grammarlens}[title={What you've built}]
One more word for telling, explaining and describing at B1.
\end{grammarlens}

\subsection*{Your turn}
\begin{itemize}
\raggedright
  \item \emph{Say it:} \emph{la smentita}
  \item \emph{Say it:} \emph{la smentita}, once more
  \item \emph{Say it:} a sentence of your own with \emph{la smentita}
  \item \emph{Recall:} say the Italian for the press conference, then the Italian for the press release, then say \emph{la smentita} again
\end{itemize}

\begin{tcolorbox}[breakable,colback=teal!4,colframe=teal!35!black,title={Before you move on}]
What does la smentita mean? (\textquotedblleft{}The denial (of a report)\textquotedblright{}.) Say it once more.
\end{tcolorbox}
